\documentclass[10pt,a4paper]{amsart}
\usepackage{amsmath,amssymb,amsthm,graphicx}
\usepackage{verbatim}
\newenvironment{reference}{\comment}{\endcomment}
%\newenvironment{reference}{}{}
\newenvironment{slogan}{\comment}{\endcomment}
\newenvironment{history}{\comment}{\endcomment}

% For commutative diagrams we use Xy-pic
\usepackage[all]{xy}

% We use 2cell for 2-commutative diagrams.
\xyoption{2cell}
% Theorem environments.
%
\theoremstyle{plain}
\newtheorem{theorem}[subsection]{Theorem}
\newtheorem{proposition}[subsection]{Proposition}
\newtheorem{lemma}[subsection]{Lemma}

\theoremstyle{definition}
\newtheorem{definition}[subsection]{Definition}
\newtheorem{example}[subsection]{Example}
\newtheorem{exercise}[subsection]{Exercise}
\newtheorem{situation}[subsection]{Situation}

\theoremstyle{remark}
\newtheorem{remark}[subsection]{Remark}
\newtheorem{remarks}[subsection]{Remarks}

\numberwithin{equation}{subsection}

% Macros
%
\def\lim{\mathop{\mathrm{lim}}\nolimits}
\def\colim{\mathop{\mathrm{colim}}\nolimits}
\def\Spec{\mathop{\mathrm{Spec}}}
\def\Hom{\mathop{\mathrm{Hom}}\nolimits}
\def\Ext{\mathop{\mathrm{Ext}}\nolimits}
\def\SheafHom{\mathop{\mathcal{H}\!\mathit{om}}\nolimits}
\def\SheafExt{\mathop{\mathcal{E}\!\mathit{xt}}\nolimits}
\def\Sch{\mathit{Sch}}
\def\Mor{\mathop{\mathrm{Mor}}\nolimits}
\def\Ob{\mathop{\mathrm{Ob}}\nolimits}
\def\Sh{\mathop{\mathit{Sh}}\nolimits}
\def\NL{\mathop{N\!L}\nolimits}
\def\CH{\mathop{\mathrm{CH}}\nolimits}
\def\proetale{{pro\text{-}\acute{e}tale}}
\def\etale{{\acute{e}tale}}
\def\QCoh{\mathit{QCoh}}
\def\Ker{\mathop{\mathrm{Ker}}}
\def\Im{\mathop{\mathrm{Im}}}
\def\Coker{\mathop{\mathrm{Coker}}}
\def\Coim{\mathop{\mathrm{Coim}}}

% Boxtimes
%
\DeclareMathSymbol{\boxtimes}{\mathbin}{AMSa}{"02}

%
% Macros for moduli stacks/spaces
%
\def\QCohstack{\mathcal{QC}\!\mathit{oh}}
\def\Cohstack{\mathcal{C}\!\mathit{oh}}
\def\Spacesstack{\mathcal{S}\!\mathit{paces}}
\def\Quotfunctor{\mathrm{Quot}}
\def\Hilbfunctor{\mathrm{Hilb}}
\def\Curvesstack{\mathcal{C}\!\mathit{urves}}
\def\Polarizedstack{\mathcal{P}\!\mathit{olarized}}
\def\Complexesstack{\mathcal{C}\!\mathit{omplexes}}
% \Pic is the operator that assigns to X its picard group, usage \Pic(X)
% \Picardstack_{X/B} denotes the Picard stack of X over B
% \Picardfunctor_{X/B} denotes the Picard functor of X over B
\def\Pic{\mathop{\mathrm{Pic}}\nolimits}
\def\Picardstack{\mathcal{P}\!\mathit{ic}}
\def\Picardfunctor{\mathrm{Pic}}
\def\Deformationcategory{\mathcal{D}\!\mathit{ef}}
\usepackage{cleveref}
\crefname{originallemma}{original lemma}{original lemmas}
\Crefname{originallemma}{Original lemma}{Original lemmas}
\makeatletter
\newlabel{divisors-lemma-characterize-affine-blowup}{{0BFL}{source}{}{sourcepointer-1}{}}
\newlabel{divisors-lemma-characterize-affine-blowup@cref}{{[originallemma][0][]0BFL}{[1][0][]source}{}{}{}}
\makeatother
\title{A sequence of point blowups with constant residue field and a fixed exceptional equation comes from a nonsingular arc}
\author{The Stacks Project Authors}
\date{stacks-project:resolve.tex:local-snapshot-20261012; GNU FDL 1.2 or later}
\begin{document}
\maketitle
\noindent\textit{Editorial note.} All delivered mathematical text is the author's own, copied from the pinned Stacks Project chapter file resolve.tex; no statement or proof step was rewritten, reordered or removed.\par\medskip
\noindent\textit{Editorial note.} The counted proof is the author's complete argument: the power estimate for the distinguished element against the powers of the maximal ideals, the rewriting of generators of the successive maximal ideals and the induction computing the lengths of the successive quotients, and the limit construction recognizing a complete discrete valuation ring together with the lifting of the map along the blowups.\par\medskip
\noindent\textit{Editorial note.} The label lemma-sequence-blowups names the counted statement, and the proof environment immediately following it in the source is delivered verbatim as its proof.\par\medskip
\noindent\textit{Editorial note.} The counted statement is carried by the author's own lemma environment in the chapter source.\par\medskip
\noindent\textit{Editorial note.} Pre-open gate: python operations/agent-authoring/tools/precard\_gate.py --source resolve.tex --label lemma-sequence-blowups --start 2802 --end 2891 --local-id 28074 returned RESULT=OK (L28074 free, occupancy FREE, public interval clear). Pre-assembly tag self-check: canonical tag resolved and not present in reference/public-items.jsonl.\par\medskip
\noindent\textit{Editorial note.} The statement and proof are copied from the pinned chapter file; the cited inputs are kept as external pointers with their canonical Stacks tags.\par\medskip
\noindent\textit{Allowed standard prerequisites (editorial).}\begin{itemize}
\item Blowing up along a regular element and the ideal of the exceptional divisor. Standard description of a blowup: the ideal of the exceptional divisor pulls back to the ideal generated by the blown up element, and multiplicities of powers translate into membership in powers of that ideal. Supplies the claim that a power of the distinguished element is not in the corresponding power of the maximal ideal of the local ring at the next step.
\item Characterization of affine blowup algebras. Divisors, the lemma characterizing an affine blowup algebra, used as one of the two ways of lifting the local homomorphism to the blowup. Supplies the explicit construction identifying the local rings of the blowups with localizations of repeated affine blowup algebras, hence the lift of the map to the limit.
\item Local rings of a blowup at a closed point and residue fields. Standard facts about the local ring of a blowup at a point and the equality of residue fields along the sequence. Supply the identification of residue fields, which allows coefficients to be chosen in the original local ring.
\end{itemize}\medskip
\ifdefined\hypertarget\hypertarget{sourcepointer-1}{}\fi
\noindent\textit{Source reference.} Original lemma 0BFL: sources/stacks-chow/tags.txt, source lines 14389--14389. Prerequisite pointer invoked at proof line(s) 2886; it names a standard fact of a sibling chapter that the counted argument uses, and every step of the counted proof built on it is retained verbatim in the delivered body.\par\medskip
\begin{lemma}
\label{lemma-sequence-blowups}
Let $X$ be a locally Noetherian scheme. Let
$$
(X, p) = (X_0, p_0) \leftarrow (X_1, p_1) \leftarrow (X_2, p_2) \leftarrow
(X_3, p_3) \leftarrow \ldots
$$
be a sequence of blowups such that
\begin{enumerate}
\item $p_i$ is closed, maps to $p_{i - 1}$, and
$\kappa(p_i) = \kappa(p_{i - 1})$,
\item there exists an $x_1 \in \mathfrak m_p$ whose image
in $\mathfrak m_{p_i}$, $i > 0$ defines the exceptional divisor
$E_i \subset X_i$.
\end{enumerate}
Then the sequence is obtained from a nonsingular arc $a : T \to X$
as above.
\end{lemma}

\begin{proof}
Let us write $\mathcal{O}_n = \mathcal{O}_{X_n, p_n}$
and $\mathcal{O} = \mathcal{O}_{X, p}$. Denote
$\mathfrak m \subset \mathcal{O}$ and $\mathfrak m_n \subset \mathcal{O}_n$
the maximal ideals.

\medskip\noindent
We claim that $x_1^t \not \in \mathfrak m_n^{t + 1}$.
Namely, if this were the case, then in the local ring
$\mathcal{O}_{n + 1}$ the element $x_1^t$ would be in the ideal of
$(t + 1)E_{n + 1}$.
This contradicts the assumption that $x_1$ defines $E_{n + 1}$.

\medskip\noindent
For every $n$ choose generators $y_{n, 1}, \ldots, y_{n, t_n}$
for $\mathfrak m_n$. As
$\mathfrak m_n \mathcal{O}_{n + 1} = x_1\mathcal{O}_{n + 1}$
by assumption (2), we can write $y_{n, i} = a_{n, i} x_1$
for some $a_{n, i} \in \mathcal{O}_{n + 1}$. Since
the map $\mathcal{O}_n \to \mathcal{O}_{n + 1}$ defines
an isomorphism on residue fields by (1) we can choose
$c_{n, i} \in \mathcal{O}_n$ having the same residue class as
$a_{n, i}$. Then we see that
$$
\mathfrak m_n = (x_1, z_{n, 1}, \ldots, z_{n, t_n}),
\quad z_{n, i} = y_{n, i} - c_{n, i} x_1
$$
and the elements $z_{n, i}$ map to elements of $\mathfrak m_{n + 1}^2$
in $\mathcal{O}_{n + 1}$.

\medskip\noindent
Let us consider
$$
J_n = \Ker(\mathcal{O} \to \mathcal{O}_n/\mathfrak m_n^{n + 1})
$$
We claim that $\mathcal{O}/J_n$ has length $n + 1$ and that
$\mathcal{O}/(x_1) + J_n$ equals the residue field. For $n = 0$
this is immediate. Assume the statement holds for $n$.
Let $f \in J_n$. Then in $\mathcal{O}_n$ we have
$$
f = a x_1^{n + 1} + x_1^n A_1(z_{n, i}) +
x_1^{n - 1} A_2(z_{n, i}) + \ldots + A_{n + 1}(z_{n, i})
$$
for some $a \in \mathcal{O}_n$ and some $A_i$ homogeneous of degree $i$
with coefficients in $\mathcal{O}_n$. Since $\mathcal{O} \to \mathcal{O}_n$
identifies residue fields, we may choose $a \in \mathcal{O}$
(argue as in the construction of $z_{n, i}$ above).
Taking the image in
$\mathcal{O}_{n + 1}$ we see that $f$ and $a x_1^{n + 1}$
have the same image modulo $\mathfrak m_{n + 1}^{n + 2}$.
Since $x_n^{n + 1} \not \in \mathfrak m_{n + 1}^{n + 2}$
it follows that $J_n/J_{n + 1}$ has length $1$ and the claim is true.

\medskip\noindent
Consider $R = \lim \mathcal{O}/J_n$. This is a quotient of
the $\mathfrak m$-adic completion of $\mathcal{O}$ hence it is
a complete Noetherian local ring. On the other hand, it is
not finite length and $x_1$ generates the maximal ideal.
Thus $R$ is a complete discrete valuation ring.
The map $\mathcal{O} \to R$ lifts to a local homomorphism
$\mathcal{O}_n \to R$ for every $n$. There are two ways to show this:
(1) for every $n$ one can use a similar procedure
to construct $\mathcal{O}_n \to R_n$ and then one can
show that $\mathcal{O} \to \mathcal{O}_n \to R_n$ factors
through an isomorphism $R \to R_n$, or (2) one can use
Divisors, Lemma \ref{divisors-lemma-characterize-affine-blowup}
to show that $\mathcal{O}_n$ is a localization of a repeated
affine blowup algebra to explicitly construct a map $\mathcal{O}_n \to R$.
Having said this it is clear that our sequence of blowups
comes from the nonsingular arc $a : T = \Spec(R) \to X$.
\end{proof}

\end{document}
